\section*{Hoofdstuk \ref{ch:g7:literal} --- Uitdrukkingen met letters}

\begin{solution}{exo:g7:literal:1}
$7x$; \quad $3y$; \quad $5ab$; \quad $x^2$; \quad $2(a + 4)$.
\end{solution}

\begin{solution}{exo:g7:literal:2}
$5ab = 5 \times a \times b$; \quad
$3(x+2) = 3 \times (x + 2)$; \quad
$x^2 + 4x = x \times x + 4 \times x$; \quad
$2a^3 = 2 \times a \times a \times a$.
\end{solution}

\begin{solution}{exo:g7:literal:3}
Voor $x = 3$: \ $5 \times 3 + 1 = 16$; \quad $3 \times 3 = 9$; \quad
$2 \times 9 - 3 = 15$; \quad $10 - 6 = 4$.
\end{solution}

\begin{solution}{exo:g7:literal:4}
$3a + 2b = 12 + 3 = 15$; \qquad
$ab + 5 = 4 \times 1.5 + 5 = 6 + 5 = 11$.
\end{solution}

\begin{solution}{exo:g7:literal:5}
$8x + 5x = 13x$; \quad
$9a - 4a + a = 6a$; \quad
$6y + 2 + 3y + 7 = 9y + 9$; \quad
$5t + 3 - 2t - 3 = 3t$.
\end{solution}

\begin{solution}{exo:g7:literal:6}
Vierkant: $P = 4z$. \omterm{def:g6:shapes:triangles}{Gelijkbenige driehoek}: $P = b + 2s$. Broden:
$p = 1.20\,n$ euro.
\end{solution}

\begin{solution}{exo:g7:literal:7}
$2n$: het dubbele; \ $n + 2$: twee meer dan $n$; \ $n^2$: het kwadraat;
\ $\frac n2$: de \omterm{def:g3:division:half}{helft}.
\end{solution}

\begin{solution}{exo:g7:literal:8}
$3(x+4) = 3x + 12$: \emph{waar} voor elke $x$ (distributiviteit --- dit is een
bewijs, testen is niet nodig).

$2x + 5 = 7x$: test $x = 2$: links $9$, rechts $14$ --- in het algemeen
\emph{niet waar} (het geldt alleen voor $x = 1$).

$x + x = x^2$: test $x = 3$: links $6$, rechts $9$ --- in het algemeen
\emph{niet waar} ($x + x = 2x$).
\end{solution}

\begin{solution}{exo:g7:literal:9}
\emph{1.} Lengte: $b + 3$. \omterm{def:g6:measure:perimeter}{Omtrek}:
\[
P = 2\bigl(b + (b + 3)\bigr) = 2(2b + 3) = 4b + 6 .
\]

\emph{2.} Voor $b = 5$: rechtstreeks zijn de zijden $5$ en $8$, dus
$P = 2 \times 13 = 26$ cm; met de formule $4 \times 5 + 6 = 26$ cm. Dat komt
overeen.
\end{solution}

\begin{solution}{exo:g7:literal:10}
De stappen met de letter $n$ gevolgd:
\[
n \ \to\ n + 6 \ \to\ 3(n + 6) = 3n + 18 \ \to\ 3n + 18 - 18 = 3n
\ \to\ \frac{3n}{n} = 3 .
\]
Iedereen krijgt $3$.
\end{solution}

\begin{solution}{exo:g7:literal:11}
Als \omterm{ex:g1:subtraction:difference}{verschil}: $A = a^2 - b^2$. Door de L in twee rechthoeken te knippen: een strook
over de volle breedte met hoogte $a - b$ en daarbovenop een smallere rechthoek:
\[
A = a(a - b) + b(a - b)
\]
(de onderste rechthoek $a$ breed en $a - b$ hoog, de bovenste $a - b$ breed en $b$
hoog geeft de variant $a(a-b) + (a-b)b$ --- hetzelfde). Voor $a = 5$ en $b = 2$:
$a^2 - b^2 = 25 - 4 = 21$, en $5 \times 3 + 2 \times 3 = 15 + 6 = 21$. Beide
uitdrukkingen komen overeen --- en dat doen ze altijd: dit is de vermomde
merkwaardige gelijkheid $a^2 - b^2 = (a + b)(a - b)$ uit
\cref{ch:g9:algebra}.
\end{solution}

\begin{solution}{pb:g7:literal:1}
\textbf{1.} De figuren $1$ tot $4$ hebben $4$, $7$, $10$ en $13$ stokjes nodig.

\textbf{2.} Een nieuw vierkant deelt één zijde met de vorige figuur, dus zijn maar
$3$ van zijn $4$ zijden nieuwe stokjes. Vanaf figuur $4$ ($13$ stokjes) kosten zes
vierkanten erbij $6 \times 3 = 18$: figuur $10$ heeft $13 + 18 = 31$ stokjes
nodig.

\textbf{3.} Lees de rij van links af: hij begint met één verticaal stokje, en elk
van de $n$ vierkanten voegt daarna een boven-, een onder- en een rechterzijde toe
--- $3$ stokjes per vierkant. Totaal: $3n + 1$; die ``$+1$'' is het allereerste
verticale stokje. Controle: $n = 1, 2, 3, 4$ geven $4, 7, 10, 13$.

\textbf{4.} $4 + 3(n - 1) = 4 + 3n - 3 = 3n + 1$: de formule van Amir werkt uit
tot precies dezelfde herleide uitdrukking. Volledige overeenstemming.

\textbf{5.} Lena: $2n + (n + 1) = 3n + 1$ --- opnieuw hetzelfde. Drie manieren van
kijken, één herleide vorm. Figuur $100$: $3 \times 100 + 1 = 301$ stokjes.

\textbf{6.} De aantallen: $3$, $5$, $7$ stokjes. Elke nieuwe driehoek leunt tegen
een bestaande zijde en voegt maar $2$ stokjes toe, vanaf $3$: figuur $n$ heeft
$3 + 2(n - 1) = 2n + 1$ stokjes nodig.

\textbf{7.} Maak $2n + 1 = 49$ ongedaan: het eerste stokje eraf laat $48$ over, en
elke driehoek staat voor $2$: $48 \div 2 = 24$ driehoeken.

\textbf{8.} In een rij van $n$ tafels biedt elke tafel haar boven- en onderzijde
($2n$ plaatsen), en de twee tafels aan de uiteinden bieden elk nog een extra zijde
($+2$): $2n + 2$ plaatsen. Controle: $4$, $6$, $8$ voor $n = 1, 2, 3$.

\textbf{9.} Eén rij: $2n + 2 \geq 30$ vraagt $2n \geq 28$, dus $14$ tafels. Deel je
dezelfde $14$ tafels in twee rijen van $7$: elke rij biedt $2 \times 7 + 2 = 16$
plaatsen, samen $32$ --- twee meer dan de enkele rij, omdat elke nieuwe rij twee
nieuwe \emph{uiteinden} meebrengt. (Elke splitsing levert $2$ plaatsen op;
cateraars weten dat.)

\textbf{10.} De L van maat $n$ heeft $n$ stippen omhoog en $n$ stippen naar de
zijkant, met de hoekstip gedeeld: $n + n - 1 = 2n - 1$ stippen. De figuren $1$,
$2$ en $3$: $1$, $3$ en $5$ stippen --- de \emph{\omterm{def:g3:numbers:evenodd}{oneven} getallen}. In elkaar
geschoven vullen de L's een vierkant van $n \times n$ stippen: het verhaal van
\cref{pb:g8:equations:1}.

\textbf{11.} $3(n + 2) - 5 = 3n + 6 - 5 = 3n + 1$: de bewering herleidt tot de
bewezen formule, dus is ze juist voor elke $n$ --- vastgesteld met algebra en niet
met voorbeelden.

\textbf{12.} Tests: $4 \geq 2$, $9 \geq 3$, $100 \geq 10$: alle drie geslaagd.
Maar $n = 0.5$ geeft $n^2 = 0.25 < 0.5$: de bewering faalt. De drie tests bewezen
niets over \emph{alle} getallen (ze vonden alleen geen problemen); dat ene
tegenvoorbeeld \emph{bewijst} dat de algemene bewering onjuist is
(\cref{met:g7:literal:test}). Een letter staat voor elk getal --- decimalen
inbegrepen.

\textbf{13.} Gebieden: $2$ punten $\to 2$; $3$ punten $\to 4$; $4$ punten
$\to 8$; $5$ punten $\to 16$. Het patroon fluistert ``elk punt verdubbelt het
aantal'': voor $6$ punten verwacht je $32$ gebieden.

\textbf{14.} Zorgvuldig tellen geeft $31$ gebieden --- geen $32$. Het
verdubbelvermoeden is dood: het doorstond vier controles en faalde bij de vijfde.
(Dit tegenvoorbeeld is beroemd genoeg om een naam te hebben, ``de valstrik van de
cirkelgebieden''; en geen sluiproutes met gelijkmatig verdeelde punten --- drie
koorden door één punt zouden gebieden laten samensmelten en de telling bederven.)

\textbf{15.} Een bewijs zou een \emph{structurele} reden nodig hebben gehad waarom
een punt erbij het aantal gebieden verdubbelt --- en die is er niet: het
verdubbelen was een toevalligheid van kleine gevallen. De luciferformules waren
veilig omdat elk van hen rechtstreeks van de constructie is afgelezen (elk vierkant
kost zichtbaar $3$ stokjes, elke tafel biedt zichtbaar $2$ plaatsen): structuur,
geen testen. Een bewezen voorspelling: $3 \times 2\,026 + 1 = 6\,079$ stokjes.
\end{solution}
